\SourcePage{GELFAND-P000455}{74}
\sout{i.e. \ldots}\qquad $x^g-x\in m_L$, i.e. $g\in G_1$.
\Unclear{GELFAND-P000455}{top-old-fragment}{The start of the crossed-out entry provisionally reads “i.e.”.}
\end{GelfandAdjacentEvidence}
\begin{GelfandAdjacentEvidence}{GELFAND-P000455}{2}

\Incomplete{GELFAND-P000455}{top-cancelled}{The remainder of the crossed-out entry is illegible.}
\end{GelfandAdjacentEvidence}
\begin{GelfandAdjacentEvidence}{GELFAND-P000455}{3}


\[
0\longrightarrow G_0/G_1\xrightarrow{\varphi}\operatorname{Gal}(\ell:k)
\]
\underline{Exercise:}\quad $\varphi$ — an isomorphism.

$2)$\quad $G_1/G_2$ — $?$
\Unclear{GELFAND-P000455}{quotient-question}{A question mark follows the dash.}
\end{GelfandAdjacentEvidence}
\begin{GelfandAdjacentEvidence}{GELFAND-P000455}{4}


\[
G_1\xrightarrow{\psi_1}\ell^*
\qquad
g\mathrel{\sim\!\longrightarrow}\overline{\tau^g/\tau}
\qquad\text{(indep. of the choice of $\tau$)}
\]
\Unclear{GELFAND-P000455}{ramification-label}{The first arrow is labelled $\psi_1$.}
\end{GelfandAdjacentEvidence}
\begin{GelfandAdjacentEvidence}{GELFAND-P000455}{5}


Its kernel is $G_2$.
\Unclear{GELFAND-P000455}{kernel-pronoun}{The word following “kernel” reads «его» (“its”).}
\end{GelfandAdjacentEvidence}
\begin{GelfandAdjacentEvidence}{GELFAND-P000455}{6}


\[
\begin{aligned}
G_i/G_{i+1}&\longrightarrow\ell^+
   \qquad\text{(addit. group of the field $\ell$).}\\
g&\xrightarrow{\psi_i}
\overline{\left(\frac{\tau^g/\tau-1}{\tau^{i-1}}\right)}
\end{aligned}
\]
\Unclear{GELFAND-P000455}{higher-label}{The arrow from $g$ is labelled $\psi_i$.}
\end{GelfandAdjacentEvidence}
\begin{GelfandAdjacentEvidence}{GELFAND-P000455}{7}

\underline{Thm.}\quad $G_{i+1}=\operatorname{Ker}\psi_i$.
\Unclear{GELFAND-P000455}{kernel-domain}{The domain on the upper line is written $G_i/G_{i+1}$. The line below has $G_{i+1}=\operatorname{Ker}\psi_i$.}
\end{GelfandAdjacentEvidence}
\begin{GelfandAdjacentEvidence}{GELFAND-P000455}{8}


$G_1$ is called the inertia subgroup of the Galois grp.

Unramified extension $\underset{\mathrm{def}}{\Longleftrightarrow} G_1=\{e\}$;

weakly ramif. $G_2=\{e\}$.
\Unclear{GELFAND-P000455}{weak-word}{The second letter is overwritten. The abbreviation reads «слаборазветвл.» (“weakly ramif.”).}
\end{GelfandAdjacentEvidence}
\begin{GelfandAdjacentEvidence}{GELFAND-P000455}{9}


purely ramified $G_0=G_1$
\Unclear{GELFAND-P000455}{last-subscript}{The last subscript is overwritten. It provisionally reads $1$.}
\end{GelfandAdjacentEvidence}


\begin{center}
\rule[.5ex]{25mm}{.35pt}\enspace$\cdot$\enspace\rule[.5ex]{35mm}{.35pt}
\end{center}


\Needspace{7\baselineskip}
\begin{GelfandAdjacentEvidence}{GELFAND-P000456}{1}
